\section{The origin completion and unconditional cotangent spanning}
\label{sec:cubic-origin-completion}

The antecedent \texttt{M2} in the abstract model's
\texttt{der} field is needed for its assertion at arbitrary
local points.  At the particular origin of the actual
weighted-power-series construction, a finite Taylor
calculation gives the required extension directly.
This section proves that comparison for the exact
algebra used in AC-G51, and removes \texttt{M2} from
the conclusion AC-G53-M1.  All arithmetic and local
existence hypotheses of AC-G51 and AC-G54 remain in
force; in particular $e\geq1$ and
$\ell\ne2,3,269$.

The exact source definitions are in
\path{Development/AutomorphicArithmetic/AutomorphicVarieties/CommutativeAlgebra/ResMapBijective.lean}.
There \texttt{WPres.part} is the image of the full
formal-series weight part under the presentation;
\texttt{WPres.A} is \texttt{gradedSub W.part}; and
\texttt{gradedSub} has underlying module
$\bigvee_j\texttt{W.part}\,j$.
Thus this algebra contains every homogeneous formal
series image and all finite sums of such images.
It is not assumed to be finitely generated over $O$.
In particular, opposite variable weights can give an
entire subring $O[[z_1z_2]]$ in weight zero.
The definition \texttt{A := WA.A} in
\path{Development/AutomorphicArithmetic/AutomorphicVarieties/UnitaryThreeGroups/ModelOfPres.lean},
\texttt{mop\_main}, is precisely this construction.

\begin{proposition}[AC-G56-jets: the exact finite jets and completion]
\label{AC-G56-jets}
Let $O$ be a complete discrete valuation ring with
uniformizer $\varpi$, and give the finitely many
variables of
\[
             P=O[[z_1,\ldots,z_d]]
                                                               \tag{G56.1}
\]
arbitrary integer weights $w_1,\ldots,w_d$.
Write $P_j$ for the formal series supported on
monomials of weight $j$, and
$P_{\rm fin}=\bigoplus_jP_j$ for their finite sums.
Let $I\subset P$ be an ideal stable under every
weight-component projection and contained in
$(z_1,\ldots,z_d)$.
Set $R=P/I$, let $A$ be the image of $P_{\rm fin}$ in
$R$, and let $x_R:R\to O$ and $x:A\to O$ be constant
term.  Then
\[
 A=\operatorname{im}(P_{\rm fin}\to R)
   =\bigoplus_j\operatorname{im}(P_j\to R).
                                                               \tag{G56.2}
\]
Put $\mathfrak m_R=(\varpi,z_1,\ldots,z_d)R$ and
$\mathfrak m_A=x^{-1}(\varpi O)$, with the same
variable notation for their images.  For every
$N\geq1$,
\[
 \begin{split}
 \mathfrak m_A&=(\varpi,z_1,\ldots,z_d)A,\\
 A\cap\mathfrak m_R^N&=\mathfrak m_A^N,\\
 A/\mathfrak m_A^N&\ \xrightarrow{\ \sim\ }\
                       R/\mathfrak m_R^N,
 \qquad
 \widehat A_{\mathfrak m_A}\ \xrightarrow{\ \sim\ }\ R .
 \end{split}                                                   \tag{G56.3}
\]

If $g_1,\ldots,g_t\in A$ are homogeneous and vanish
at $x$, put $J=(g_1,\ldots,g_t)A$ and
$S=R/(g_1,\ldots,g_t)R$.  Then the natural map
$A/J\to S$ is injective, and, for the maximal ideal
$\mathfrak n$ induced by $x\bmod\varpi$,
\[
             \widehat{(A/J)}_{\mathfrak n}\simeq S .
                                                               \tag{G56.4}
\]
These are the actual completions at the specified
closed point; no degree-zero finiteness hypothesis
is used.
\end{proposition}
\begin{proof}
A finite sum of homogeneous series maps to a finite
sum of homogeneous images.  Conversely, each such
image has a homogeneous lift by definition, so a
finite sum has a lift in $P_{\rm fin}$.
If a finite sum of different homogeneous series
maps to zero, its total lift lies in $I$; stability
of $I$ under the weight projections puts each
summand in $I$.  This proves the directness in
(G56.2) and the asserted surjection
$P_{\rm fin}\twoheadrightarrow A$.
It also proves directly that this $A$ is the source's
\texttt{WPres.A}.

For every $F\in P_{\rm fin}$, assign each nonconstant
monomial to its first variable with positive
exponent.  Dividing the corresponding terms by that
variable gives
\[
                F=F(0)+\sum_{i=1}^d z_iG_i,
                   \qquad G_i\in P_{\rm fin}.
                                                               \tag{G56.5}
\]
Indeed each $G_i$ has weights in the finite weight
support of $F$ shifted by $-w_i$.
If $x(F)\in\varpi O$, (G56.5) proves the first
identity of (G56.3); the reverse containment follows
by applying $x$.  This ideal is maximal since its
quotient is $O/\varpi O$.

We record the elementary homogeneous-ideal
contraction used for the powers and for (G56.4).
If $h_1,\ldots,h_v\in A$ are homogeneous of weights
$c_1,\ldots,c_v$, then
\[
          A\cap(h_1,\ldots,h_v)R=(h_1,\ldots,h_v)A .
                                                               \tag{G56.6}
\]
Choose homogeneous lifts $H_i\in P_{c_i}$.
Given $a\in A$ in the left side, choose a finite-weight
lift $F$ of $a$ and arbitrary lifts $B_i\in P$ of
coefficients expressing $a=\sum_i b_ih_i$.
Then $F-\sum_iB_iH_i\in I$.
For each weight $j$, its weight component gives in
$R$
\[
       \overline{F_j}
          =\sum_i\overline{(B_i)_{j-c_i}}\,h_i .
                                                               \tag{G56.7}
\]
Every coefficient on the right belongs to $A$.
Sum (G56.7) over the finite weight support of $F$.
This proves one inclusion of (G56.6), and the other
is immediate.  This is the full component argument
of the source
\path{Development/AutomorphicArithmetic/AutomorphicVarieties/CommutativeAlgebra/HomogSpan.lean},
\texttt{homog\_span}.

Apply (G56.6) to the finite family of monomials of
total degree $N$ in $\varpi,z_1,\ldots,z_d$.
They are homogeneous, with $\varpi$ of weight zero,
and generate $\mathfrak m_R^N$ and
$\mathfrak m_A^N$ in their respective rings.
This proves the contraction identity in (G56.3).
Every formal series is congruent modulo
$(z_1,\ldots,z_d)^N$ to its finite total-degree
truncation, which belongs to $P_{\rm fin}$.
Consequently $A\to R/\mathfrak m_R^N$ is surjective.
Its kernel is the contraction just proved, giving
the finite-jet isomorphism.
The ring $R$ is complete for $\mathfrak m_R$:
$P$ is a complete Noetherian local ring and its
ideals are closed.  Taking inverse limits proves
the completion assertion.

For the last statement, (G56.6) with $h_i=g_i$
gives injectivity of $A/J\to S$.
The preimage of $(g_i)R$ in $P$ is stable under
weight projections: the same calculation as
(G56.7), applied to an arbitrary series in that
preimage and to each weight separately, proves
this assertion.  It is contained in the
constant-term kernel because every $g_i$ vanishes
at the origin.  The image of $P_{\rm fin}$ in
$S$ is exactly $A/J$.
Apply the already proved finite-jet and completion
statements to this new quotient of $P$.
This proves (G56.4) without an appeal to flatness
of a completion or to finite generation over $O$.
\end{proof}

\begin{lemma}[AC-G56-der: every origin derivation extends]
\label{AC-G56-der}
In the setting of AC-G56-jets, let $Q$ be any
$O$-module and let $D:A\to Q$ be an $O$-linear
derivation at $x$.
There is a unique $O$-linear derivation
$\widetilde D:R\to Q$ at $x_R$ restricting to $D$.
For $F\in P$ it is given by the finite formula
\[
       \widetilde D(\overline F)
           =\sum_{i=1}^d [z_i]F\,D(\overline z_i).
                                                               \tag{G56.8}
\]
If $D(J)=0$, it descends to a derivation
$D_S:S\to Q$ at the origin of $S$.
If $D$ has one weight $b$, its extension has that
same weight.

For $Q=\overline{\mathbb Q}_\ell$ and the actual
coefficient ring $O$, $D_S$ takes values in a
finite coefficient extension and is continuous.
Thus
\[
             s\longmapsto x_0(s)+\varepsilon D_S(s)
                                                               \tag{G56.9}
\]
is an actual dual-number point of $S$ with
first component $\overline{x}_0$.
No \texttt{M2} hypothesis is required.
\end{lemma}
\begin{proof}
Pull $D$ back along the proved surjection
$P_{\rm fin}\to A$.
For $F\in P_{\rm fin}$, split off its constant
and degree-one terms.  Assign each remaining
monomial to its first two variable factors,
allowing the same variable twice.  This gives
\[
 F=F(0)+\sum_i[z_i]F\,z_i+
                 \sum_{i,k}z_i z_k H_{ik},
       \qquad H_{ik}\in P_{\rm fin}.
                                                               \tag{G56.10}
\]
There are only finitely many pairs $(i,k)$; the
weight support of $H_{ik}$ is contained in the
finite support of $F$ shifted by $-w_i-w_k$.
A derivation at constant term kills constants
and every product of two origin-vanishing
elements.  Therefore (G56.10) proves (G56.8)
on $P_{\rm fin}$, algebraically and without a
continuity assumption.

Define the right side of (G56.8) on all of $P$.
It is $O$-linear, and the coefficient identity
\[
     [z_i](FG)=F(0)[z_i]G+G(0)[z_i]F
                                                               \tag{G56.11}
\]
proves its derivation rule at constant term.
Let $F\in I$.  For each weight $j$ among the
finitely many $w_i$, the component $F_j$
belongs to $I\cap P_{\rm fin}$.
The derivative of this component is zero because
the pulled-back $D$ kills the kernel of
$P_{\rm fin}\to A$.
Grouping the finite sum (G56.8) by the values
of $w_i$ expresses it as the sum of these
zero component derivatives.  Thus it kills all
of $I$, not merely its finite-weight elements,
and descends to $R$.

It agrees with $D$ because every element of $A$
has a finite-weight lift.  Conversely, any
derivation on $R$ pulls back to a derivation on
$P$ at constant term.  Decomposition (G56.10)
still holds with arbitrary formal series
$H_{ik}\in P$, so the same calculation forces
(G56.8).  This proves uniqueness without
assuming continuity.
If $D(J)=0$, the extension kills each $g_i$.
Since $x_R(g_i)=0$, its product rule then
kills every $r g_i$ for $r\in R$, and hence the
full ideal $(g_i)R$.  It therefore descends to
$S$.

If $D$ has weight $b$, then
$D(\overline z_i)=0$ whenever $w_i\ne b$.
Formula (G56.8) consequently vanishes on every
homogeneous part of other weight.  This proves
the weight assertion.

Finally choose a finite extension $L/\mathbb Q_\ell$
containing $\operatorname{Frac}(O)$ and the finitely
many values $D(\overline z_i)$.
The formula takes values in $L$.
For $N\geq1$, a series in
$(\varpi,z_1,\ldots,z_d)^N$ has constant term
in $\varpi^N O$ and each linear coefficient in
$\varpi^{N-1}O$.
If $C=\max(1,|D(\overline z_1)|_\ell,\ldots,
|D(\overline z_d)|_\ell)$, it follows
that
\[
       |D_S(s)|_\ell\leq C|\varpi|_\ell^{N-1}
                \quad(s\in\mathfrak m_S^N).
                                                               \tag{G56.12}
\]
To apply this estimate, lift such an $s$ by an
element of the indicated power in $P$; the
maximal ideal of $S$ is the image of the
maximal ideal of $P$.  This proves continuity.
The first component $x_0$ is the actual origin
and sends $\mathfrak m_S$ into $\varpi O$, so
it satisfies the source's \texttt{IsLocPt}.
The derivation rule is exactly the multiplicativity
of (G56.9).  This is the required actual dual point.
\end{proof}

\begin{theorem}[AC-G56: the actual cubic model satisfies \texttt{M1}]
\label{AC-G56}
Let $\mathcal M_j$ be the actual model constructed
in AC-G51, with its local construction extended
by AC-G54 where necessary.
Under their retained arithmetic hypotheses,
including \texttt{CohInp}, \texttt{GD} and the
full \texttt{LemmaP}, the source condition
$\texttt{M1}(\mathcal M_j.Q)$ holds for every
coefficient prime in their range.
In particular the equation families of weights
$-1,-2,0$ satisfy the source's exact saturated
cotangent spanning conditions
\texttt{SpansCot}, without assuming \texttt{M2}.
\end{theorem}
\begin{proof}
The ring in (G51.7), its homogeneous kernel,
and the origin satisfy AC-G56-jets.
The model algebra is exactly its $A$, by the
literal definitions recalled above; its ideal
$J$ is generated by the finitely many homogeneous
global and normalization equations of (G51.4).
The equalities (G51.8)--(G51.9) identify its
completed quotient with the actual normalized
ring $S$.  Thus AC-G56-der applies to the actual
model, not to an auxiliary polynomial subalgebra.

Let $b\in\{-2,-1,0\}$ and let $D:A\to
\overline{\mathbb Q}_\ell$ be a weight-$b$
derivation at $x$ killing $J$.
Lemma AC-G56-der gives an actual dual point
$F$ on $S$ above $x_0$ with derivative extending
$D$.
Each universal entry of weight different from
$b$ has a homogeneous preimage of that weight
in $A$, by AC-G51.
Its derivative is therefore zero.
Multiplication by the inverse of the diagonal
basepoint matrix multiplies each entry by a
unit and preserves its matrix weight, so the
tangent cocycle of $F$ has weight $b$.
The zero conclusion of AC-G53 gives
$D_S=0$, hence $D=0$.
This is exactly the step for which AC-G53-M1
previously invoked the conditional
\texttt{Model.der} field; here its antecedent
\texttt{M2} has been replaced by the proved
origin extension.

For completeness, the cotangent implication uses
the precise saturated definition, not an integral
generation assertion.  Put $I_x=\ker x$ and
\[
 I_x^{[2]}=\{a\in A:\text{some }0\ne c\in O
                              \text{ has }ca\in I_x^2\}.
                                                               \tag{G56.13}
\]
It is graded because $I_x$ is graded and weight
projection preserves $I_x^2$.
If the equations of weight $b$ failed to span
the weight-$b$ part of
$(I_x/I_x^{[2]})\otimes_O\operatorname{Frac}(O)$,
there would be a nonzero linear functional on
that part vanishing on those equations.
Extend it by zero on the other weights.
Composing it with
$a\mapsto a-x(a)$ defines an $O$-linear
weight-$b$ derivation: the error in the product
rule lies in $I_x^2$.
It kills every equation of other weight and
the specified weight-$b$ equations.  Because
all equations vanish at $x$, the derivation
rule then makes it zero on their entire ideal
$J$.  The preceding paragraph forces this
derivation to be zero, a contradiction.
The spanning assertion over
$\operatorname{Frac}(O)$ means that for each
cotangent element some nonzero $c\in O$
makes it an $O$-linear combination of the
equation classes; this is exactly
\texttt{SpansCot}.

The rank inequality $r\geq1$ and the required
\texttt{SmoothPt} are already part of AC-G51.
Together with these three spanning statements
they give the full literal \texttt{M1}.
This last implication is also the algebraic
\texttt{spans\_of\_der} argument in
\path{Development/AutomorphicArithmetic/AutomorphicVarieties/UnitaryThreeGroups/GradedModelSpansCotOfRigidWeights.lean};
its derivation-extension input has now been
proved at the origin directly.
\end{proof}

\paragraph{Scope of the improvement.}
This proves \texttt{M1} for the actual cubic model.
It does not assert the arbitrary-point extension
without \texttt{M2}, does not prove \texttt{M2},
and supplies no numerical cotangent length upper
bound.  The actual Type I automorphy and projected
mixed tangent requirements remain separate.
The \texttt{M3} implication in AC-G55 still has
its stated \texttt{M2} antecedent.

\paragraph{Geometric dictionary.}
The finite-jet quotients in (G56.3) describe
infinitesimal neighborhoods of the specified
closed point, and (G56.4) identifies their
formal limit with the actual deformation ring.
The lisse finite-cover objects being deformed
retain the already audited SGA~1, Expos\'e~V,
\S\S4 and~7 dictionary.  Neither that dictionary
nor a new archival claim is used in place of
the explicit completion and tangent proofs above.
